\documentclass[11pt]{article}
\usepackage[utf8]{inputenc}
\usepackage[T1]{fontenc}
\usepackage{amsmath, amssymb, amsthm, mathtools}
\usepackage{geometry}
\usepackage{array}
\usepackage{booktabs}
\usepackage{enumitem}
\usepackage{microtype}
\usepackage{xurl}
\usepackage[hidelinks]{hyperref}

\geometry{margin=1in}

\title{The General Theory of Finite Obstruction Calculus}
\author{Jeremy H. Carroll}
\date{May 2026}

\newtheorem{theorem}{Theorem}[section]
\newtheorem{lemma}[theorem]{Lemma}
\theoremstyle{definition}
\newtheorem{definition}[theorem]{Definition}
\newtheorem{axiom}[theorem]{Axiom}

\begin{document}

\maketitle

\section*{Publication Main}

\begin{quote}
All files are available open source at:\par\smallskip\noindent\url{https://github.com/dbccpage/general_theory_of_finite_obstruction_calculus}
\end{quote}

\begin{quote}
This file is the publication-facing spine extracted from the working tome. Full proofs are in \path{2_GTFOC_Proof_Archive.md}.\par\smallskip
LLM, Kaggle, domain, quantum, and miscellaneous workbench material is kept in \path{4_GTFOC_Tertiary_Workbench.md}.
\end{quote}

\section*{Reader's Guide}

\textbf{Finite Obstruction Calculus (FOC)} is the finite mathematical engine that evaluates a certified regime by representing defects in a finite cochain complex, quotienting out exact repairs, measuring residual obstruction, and classifying the result as exact, survivor, open, lift-required, refused, or otherwise terminalized. The main paper keeps only the finite obstruction core: realization packet, cochain complex, \(Q^1\), \(\Phi\), \(H^1\), closure trichotomy, observer passivity, context-witness lift trigger, minimal resolver, and no-silent-mutation laws. The proof archive preserves the full proof text. The tertiary workbench keeps applications, implementation ledgers, toy models, speculative extensions, and external bridges out of the publication spine.

\begin{abstract}
We formalize a finite obstruction-repair calculus for contextual symbolic systems. A level is a finite tuple

\[
\mathcal M=(\mathcal T,\mathcal C,\rho,C^\bullet,\omega,p)
\]

where \(\mathcal T\) and \(\mathcal C\) are finite categories, \(\rho:\mathcal C\to\mathrm{End}(\mathcal T)\) is a context action, and

\[
C^0\xrightarrow{d_0}C^1\xrightarrow{d_1}C^2,
\qquad d_1d_0=0
\]

is a finite cochain complex. Exact repairs form \(B^1=\operatorname{im}(d_0)\); quotient obstruction is

\[
Q^1=C^1/B^1;
\]

and closed residual obstruction is

\[
H^1=\ker(\bar d_1)\cong\ker(d_1)/\operatorname{im}(d_0).
\]

A declared finite gauge gives a quotient obstruction magnitude \(\Phi\). Context-visible non-congruence supplies the lift trigger. Licensed lifts require explicit witness, resolver, minimality, preservation, and no-silent-mutation certificates.
\end{abstract}

\section*{Non-Claims}

This theory does not claim that every truth is finite. It does not claim that every obstruction is physical, cognitive, quantum, analytic, or metaphysical. It does not claim that a positive obstruction mass automatically licenses a lift. It does not claim Langlands, zeta, quantum-information, continuum, consciousness, or empirical-world consequences without a separate finite descent/certificate path.

The narrow claim is:

\textbf{Narrow claim:} finite obstruction authority requires finite realization, quotient obstruction, certificates, and terminal/refusal semantics.

\section{Synchronization with Finite Authority Systems}

The finite obstruction calculus is the local verifier inside the Finite Authority System arena. In the global authority pipeline,

\[
\mathcal C_0
\xrightarrow{\mathcal R}
\mathcal C
\xrightarrow{\rho}
\mathcal S
\xrightarrow{D}
\mathcal P
\xrightarrow{V}
\mathbb V
\xrightarrow{\tau}
\mathbb T,
\]

this monograph supplies the internal logic of the transition

\[
D\to\mathbb V\to\mathbb T.
\]

No claim in this calculus has theorem-grade or authority-grade status unless it is coupled to a finite realization packet and replayable certificate.

\section{Primitive Distinguishability}

\begin{axiom}[T0]
Distinguishability is primitive. Sameness is certified zero-difference inside a declared difference regime.
\end{axiom}

For a declared difference detector \(\delta\),

\[
a=b \iff \delta(a,b)=0.
\]

The zero has meaning only because the difference detector exists. Thus exactness is not absence of structure; it is certified absorption by a declared repair image.

\section{Finite Realization Gate}

A theorem-grade claim \(X\) in the finite obstruction calculus requires a declared finite realization packet

\[
\mathcal G=(K,\mathcal C,\rho,C^\bullet,E,\mathcal O,V,\tau,\mathcal K,\Lambda)
\]

where

\[
C^\bullet:C^0\xrightarrow{d_0}C^1\xrightarrow{d_1}C^2,
\qquad d_1d_0=0.
\]

The gate theorem is:

\[
\boxed{
\mathsf{TheoremGrade}(X)\Rightarrow
K\wedge(\mathcal C,\rho)\wedge C^\bullet\wedge E\wedge\mathcal O\wedge V\wedge\tau\wedge\mathcal K\wedge\Lambda.
}
\]

Equivalently, absence of any required component gives a typed realization/refusal terminal, not a proof of falsity.

\section{Ambient Finite Core}

Let

\[
\mathsf{FinCat},\qquad \mathsf{FinSet},\qquad \mathsf{Vect}^{fd}_{\mathbb R}
\]

denote finite small categories, finite sets, and finite-dimensional real vector spaces.

For a finite category \(\mathcal C\),

\[
\widehat{\mathcal C}_{fin}=\mathsf{FinSet}^{\mathcal C^{op}}
\]

is the finite presheaf environment. If a finite Grothendieck topology \(J\) is explicitly supplied, one may form

\[
\mathsf{Sh}(\mathcal C,J)\subseteq\widehat{\mathcal C}_{fin}.
\]

No sheaf topology is automatic.

\section{MSS-N Level}

An MSS level is a finite tuple

\[
\mathcal M_n=(\mathcal T_n,\mathcal C_n,\rho_n,L_n,D_n,O_n,\delta_n,\chi_n,\omega_n,\mu_n)
\]

where \(\mathcal T_n\) is a finite term/state category, \(\mathcal C_n\) is a finite admissible context category, \(\rho_n:\mathcal C_n\to\mathsf{End}(\mathcal T_n)\) is a context action, \(L_n\) is the repair generator set, \(D_n\) is the distinction-coordinate set, \(O_n\) is the closure/obstruction-coordinate set, and \(\omega_n\) is a declared positive gauge.

The cochain carrier is

\[
C^0_n=\mathbb R^{L_n},\qquad
C^1_n=\mathbb R^{D_n},\qquad
C^2_n=\mathbb R^{O_n},
\]

with

\[
d_{0,n}=\delta_n,
\qquad
d_{1,n}=\chi_n,
\qquad
d_{1,n}d_{0,n}=0.
\]

\section{Defects, Exact Repair, and Quotient Obstruction}

For a submitted finite state \(s\), an edge defect has the form

\[
\delta_e(s)=r_e(s(v))-s(u).
\]

The exact repair subspace is

\[
B^1_n=\operatorname{im}(d_{0,n})\subseteq C^1_n.
\]

The quotient obstruction space is

\[
Q^1_n=C^1_n/B^1_n,
\qquad
\pi_n:C^1_n\to Q^1_n.
\]

Because \(d_{1,n}d_{0,n}=0\), \(d_{1,n}\) descends to

\[
\bar d_{1,n}:Q^1_n\to C^2_n,
\qquad
\bar d_{1,n}\pi_n=d_{1,n}.
\]

The closed quotient sector is

\[
H^1_n=\ker(\bar d_{1,n})\cong\ker(d_{1,n})/\operatorname{im}(d_{0,n}).
\]

Elements of \(H^1_n\setminus\{0\}\) are closed non-exact residual classes.

\section{Normed Cokernel and Obstruction Magnitude}

A positive weight system gives

\[
\|x\|_{1,\omega_n}=\sum_{d\in D_n}\omega_{n,d}|x_d|.
\]

The quotient obstruction norm is

\[
\Phi_n([x])=\inf_{\alpha\in C^0_n}\|x-d_{0,n}\alpha\|_{1,\omega_n}.
\]

The closure mass is

\[
\Gamma_n([x])=\|\bar d_{1,n}([x])\|.
\]

\paragraph{Safe norm principle:} independent finite defects add by outer \(\ell^1\) after local gauges are declared.

Local gauges may be scalar \(\ell^1\), vector \(\ell^2\), trace norm, or another declared finite gauge. The core does not assert that \(\ell^1\) is universal.

\section{Main Theorems}

\begin{theorem}[T1 — Finite Realization Gate]
No theorem-grade finite obstruction claim exists without a finite realization packet
\[
(K,\mathcal C,\rho,C^\bullet,E,\mathcal O,V,\tau,\mathcal K,\Lambda).
\]
The failed claim is not false by default; it terminalizes as a missing-carrier, missing-context, missing-complex, missing-operator, missing-evaluator, missing-certificate, or missing-refusal-semantics condition.
\end{theorem}

\begin{theorem}[T2 — Normed Cokernel and Quotient Obstruction]
Let
\[
C^0\xrightarrow{d_0}C^1\xrightarrow{d_1}C^2
\]
be a finite-dimensional real cochain complex with \(d_1d_0=0\). Let \(C^1\) carry a declared norm. Define
\[
B^1=\operatorname{im}(d_0),
\qquad
Q^1=C^1/B^1,
\]
and
\[
\Phi([x])=\inf_{\alpha\in C^0}\|x-d_0\alpha\|.
\]
Then \(Q^1\) is the cokernel of \(d_0\), \(\Phi\) is the quotient norm on \(Q^1\), the infimum is attained, \(d_1\) descends to \(\bar d_1:Q^1\to C^2\), and
\[
H^1=\ker(\bar d_1)\cong\ker(d_1)/\operatorname{im}(d_0).
\]
\end{theorem}

\begin{theorem}[T3 — Primal-Dual Certificate Identity]
For \(d_0:C^0\to C^1\cong\mathbb R^n\) and positive weights \(\omega_i\), define
\[
\Phi([x])=
\min_{\alpha\in C^0}\sum_i\omega_i|x_i-(d_0\alpha)_i|.
\]
Then
\[
\boxed{
\Phi([x])=
\max_{\varphi\in(C^1)^*}
\{\varphi(x): d_0^*\varphi=0,\ |\varphi_i|\le \omega_i\}.
}
\]
The primal minimum and dual maximum are attained. A replayable certificate consists of primal repair data, dual annihilator data, and a primal-dual equality check.
\end{theorem}

\begin{theorem}[T4 — Context Witness Lift Trigger]
Let \(S=(\mathcal T,\mathcal C,\rho;C^\bullet,\omega,p)\) with
\[
p:\operatorname{Ob}(\mathcal T)\to Q^1.
\]
Define
\[
x\sim y\iff p(x)=p(y).
\]
A lift is triggered only by context-visible non-congruence:
\[
\boxed{
\mathsf{LiftTrigger}(S)
\iff
\exists x,y,c:
 x\sim y\ \wedge\ \rho(c)x\not\sim\rho(c)y.
}
\]
Obstruction mass alone does not imply licensed lift:
\[
\Phi(q)>0\not\Rightarrow\mathsf{LicensedLift}.
\]
\end{theorem}

\begin{theorem}[T5 — Minimal Resolver Existence]
Let \(w\) be a typed context witness, and let \(\mathcal G(w)\) be a declared finite candidate resolver class. If valid resolvers are decidable, nonempty, and ordered by a finite well-founded minimality order, then there exists a minimal valid resolver
\[
G^*\in\mathcal G^{\mathrm{vRes}}(w).
\]
The theorem proves existence, not uniqueness. If incomparable minimal resolvers exist and no declared equivalence identifies them, the valid terminal is nonunique-minimality or refusal, not silent choice.
\end{theorem}

\begin{theorem}[T6 — Generator Growth Law]
If finite generators update by
\[
G_{n+1}=G_n\sqcup\coprod_{w\in W_n}\Delta_G(w),
\]
with finite witness set \(W_n\), then
\[
\boxed{
S(n+1)=S(n)+\sum_{w\in W_n}\Delta(w).
}
\]
\end{theorem}

\begin{theorem}[T7 — Skeleton Preservation / No Silent Mutation]
Let \(F_n:\mathcal M_n\to\mathcal M_{n+1}\) be a licensed lift, and let \(U_n:\mathcal M_{n+1}\to\mathcal M_n\) be its forgetful comparison. Define
\[
\mathsf{LF}_n=U_nF_n.
\]
For certified old skeleton data \(q\in\operatorname{Skel}_n\), preservation requires
\[
\boxed{
q\in\operatorname{Skel}_n\Rightarrow U_nF_n(q)=q.
}
\]
If \(U_nF_n(q)\ne q\) without an explicit terminal tag declaring refinement, quotienting, invalidation, or refusal, then the lift commits silent mutation.
\end{theorem}

\section{Closure Trichotomy}

For \(q\in Q^1_n\), define
\[
\mathsf{Exact}_n(q)\iff\Phi_n(q)=0,
\]
\[
\mathsf{Survivor}_n(q)\iff\Phi_n(q)>0\wedge\bar d_{1,n}(q)=0,
\]
\[
\mathsf{Open}_n(q)\iff\Phi_n(q)>0\wedge\bar d_{1,n}(q)\ne0.
\]
Finite dimensionality gives the disjoint partition
\[
Q^1_n=\mathsf{Exact}_n\sqcup\mathsf{Survivor}_n\sqcup\mathsf{Open}_n.
\]
A survivor is a closed non-exact quotient residual. An open obstruction is not yet structure; it requires repair, lift, refusal, or budget boundary.

\section{Observers and Epistemic Quotients}

An observer is a finite, deterministic, repair-compatible quotient map
\[
\mathcal O:Q^1_n\to Q^1_{\mathcal O}.
\]
Observer quotients are epistemic structure. They erase or identify distinctions. They do not create obstruction.

Cross-observer stability is always relative to a declared finite observer family
\[
\mathfrak O_n=\{\mathcal O_1,\dots,\mathcal O_k\}.
\]
A class is observer-invariant only relative to a proven observer class.

\section{Terminal Semantics}

A verified obstruction class terminalizes into one of the following schematic statuses:

\begin{center}
\begin{tabular}{ll}
\toprule
Terminal & Meaning \\
\midrule
\(\mathsf{Exact}\) & zero obstruction under declared repair image \\
\(\mathsf{Survivor}\) & closed non-exact residual \\
\(\mathsf{Open}\) & nonzero obstruction with nonzero closure mass \\
\(\mathsf{LicensedLift}\) & context witness plus valid lift contract \\
\(\mathsf{BudgetUnknown}\) & finite budget exhausted without authority \\
\(\mathsf{NoRule}\) & no declared rule licenses the move \\
\(\mathsf{TypedRefusal}\) & claim falls outside the current authority regime \\
\bottomrule
\end{tabular}
\end{center}

Non-exact obstruction never silently authorizes exactness.

\section{Compact Formula Sheet}

\[
C^0\xrightarrow{d_0}C^1\xrightarrow{d_1}C^2,
\qquad d_1d_0=0.
\]
\[
B^1=\operatorname{im}(d_0),
\qquad
Q^1=C^1/B^1,
\qquad
\pi:C^1\to Q^1.
\]
\[
\bar d_1([x])=d_1x.
\]
\[
H^1=\ker(\bar d_1)\cong\ker(d_1)/\operatorname{im}(d_0).
\]
\[
\Phi([x])=\inf_{\alpha\in C^0}\|x-d_0\alpha\|.
\]
\[
\Gamma([x])=\|\bar d_1([x])\|.
\]
\[
x\sim y\iff p(x)=p(y).
\]

\textbf{Lift trigger:} there exist finite objects \(x,y\) and a context \(c\) such that \(x\sim y\) but \(\rho(c)x\not\sim\rho(c)y\).

\section{Minimal Theorem Statement}

A finite contextual obstruction system has theorem-grade authority only after finite realization. Its exact repair classes form \(B^1=\operatorname{im}(d_0)\), its obstruction quotient is \(Q^1=C^1/B^1\), and its closed residual sector is \(H^1=\ker(\bar d_1)\). A positive obstruction class blocks exactness but does not by itself license lift. Lift requires context-visible non-congruence plus a finite resolver/lift contract. Any representation extension that changes certified skeleton data without explicit terminal tagging commits silent mutation.

\section{Safe Claims}

\begin{enumerate}
    \item The finite obstruction core is a finite cochain complex with certified quotient obstruction.
    \item Exactness is zero class in \(Q^1\).
    \item Closed non-exact residuals are nonzero elements of \(H^1\).
    \item \(\Phi\) is a declared finite-gauge quotient norm, not a universal metaphysical norm.
    \item Context-visible non-congruence is the core lift trigger.
    \item Licensed lift requires witness, resolver, minimality, preservation, and replay certificates.
    \item Observer quotients do not create obstruction.
    \item Speculative bridges require separate finite descent packets.
\end{enumerate}

\section{Boundary}

The finite obstruction calculus is expansive by extension, not by implication. New domains may be represented, reflected, or connected only after a finite realization and certificate path is supplied.

\end{document}
